% capitolo 2 (dopo introduzione): Analisi delle serie storiche finanziarie
\chapter{Analisi delle serie storiche finanziarie}
\label{cap:analisi-serie}

Per comprendere le dinamiche dei mercati è necessario condurre un'analisi delle serie storiche, il cui scopo è tradurre le fluttuazioni dei 
prezzi in metriche misurabili di rendimento e volatilità. Questo capitolo introduce tali concetti, indispensabili per poter valutare l'impatto dei dazi.

Il capitolo analizza le modalità di calcolo dei rendimenti e definisce il concetto di volatilità, per poi concentrarsi sui principali 
comportamenti tipici (i cosiddetti «fatti stilizzati») che caratterizzano i mercati finanziari.

Lo studio di queste dinamiche serve a mostrare che i modelli lineari classici non sono sufficienti per descrivere la realtà dei mercati, 
rendendo necessario l'utilizzo dei modelli econometrici della famiglia GARCH, che vengono affrontati nei capitoli successivi.


\section{I rendimenti finanziari}
Nei mercati finanziari, il prezzo indica la quotazione puntuale di un asset, mentre il rendimento ne misura la variazione relativa nel tempo\cite{libro:investments}. 
Nell'analisi delle serie storiche, l'oggetto di studio primario non è rappresentato dai prezzi degli asset, bensì dai loro rendimenti. 
Questa scelta metodologica è motivata per due ragioni\cite{campbell:econometrics}.

La prima ragione è che il rendimento fornisce un riassunto completo dell'opportunità di investimento. Essendo una misura \gls{scale-free}, 
permette di confrontare in modo diretto strumenti finanziari differenti. 
Inoltre, le serie storiche dei rendimenti presentano proprietà statistiche favorevoli, prima tra tutte la \gls{stazionarieta}, che agevola 
l'identificazione di modelli predittivi e lo studio della volatilità.

Per misurare la variazione del prezzo nel tempo si distinguono due tipi di rendimento: il rendimento semplice uniperiodale e il 
rendimento composto a tempo continuo (log-rendimento).


\subsection{Rendimento semplice uniperiodale}
Sia $P_t$ il prezzo di un asset alla data $t$ e $P_{t-1}$ il suo prezzo nel periodo precedente, e si assuma che l'asset non paghi dividendi\cite{campbell:econometrics}.
Il rendimento netto semplice uniperiodale $R_t$, nell'intervallo di tempo $[t - 1, t]$, è definito come:
\[
    R_t = \frac{P_t}{P_{t-1}} - 1 = \frac{P_t - P_{t-1}}{P_{t-1}}
\]
e di conseguenza il rendimento lordo semplice, cioè il capitale totale dopo la detenzione dell'asset per un periodo è pari a $1 + R_t = \frac{P_t}{P_{t-1}}$.

Partendo da questo è possibile estendere la misurazione su orizzonti temporali più ampi, calcolando il rendimento lordo di un asset sugli ultimi $k$ periodi.
Denotato come $1 + R_t(k)$, equivale al prodotto dei $k$ rendimenti uniperiodali, ovvero:
\begin{equation*}
\begin{split}
    1 + R_t(k) &= (1 + R_t) \cdot (1 + R_{t-1}) \cdots (1 + R_{t-k+1}) \\
               &= \frac{P_t}{P_{t-1}} \cdot \frac{P_{t-1}}{P_{t-2}} \cdots \frac{P_{t-k+1}}{P_{t-k}} \\
               &= \frac{P_t}{P_{t-k}}
\end{split}
\end{equation*}
Il corrispondente rendimento netto multiperiodale si ottiene sottraendo l'unità: $R_t(k) = \frac{P_t}{P_{t-k}} - 1$.

Questo mostra che il rendimento semplice non gode della proprietà di additività temporale, rendendo necessaria la moltiplicazione. Per questo motivo si preferisce utilizzare i log-rendimenti.


\subsection{Rendimento composto a tempo continuo}
Un'alternativa al rendimento semplice uniperiodale è rappresentata dal rendimento composto a tempo continuo, noto anche come log-rendimento.

Il rendimento composto a tempo continuo calcolato su un singolo periodo $r_t$ è definito come il logaritmo naturale del rendimento lordo semplice:
\[
    r_t = \ln(1 + R_t) = \ln\left(\frac{P_t}{P_{t-1}}\right)
\]

Definendo $p_t = \ln(P_t)$, attraverso le proprietà dei logaritmi la formula si semplifica in: 
\[
    r_t = p_t - p_{t-1}
\]

Il vantaggio emerge nel calcolo su orizzonti temporali ampi, in quanto il log-rendimento multiperiodale $r_t(k)$ è calcolabile come somma dei $k$ log-rendimenti:
\begin{align*}
    r_t(k) &= \ln(1 + R_t(k)) \\
           &= \ln(1 + R_t) + \ln(1 + R_{t-1}) + \dots + \ln(1 + R_{t-k+1}) \\
           &= r_t + r_{t-1} + \dots + r_{t-k+1}
\end{align*}

I rendimenti composti a tempo continuo multiperiodali corrispondono alla somma dei rendimenti composti a tempo continuo uniperiodali e ciò li rende preferibili per l'analisi econometrica e lo studio della volatilità.


\section{La volatilità}
\label{sec:volatilita}
In ambito finanziario la volatilità non coincide con il rischio, ma se interpretata come incertezza diventa un elemento chiave per le decisioni 
di investimento e per la creazione dei portafogli\cite{poon:forecasting}. Infatti, il punto di partenza per la valutazione del rischio 
è una corretta previsione della futura volatilità dei prezzi degli asset.

Per via del forte legame tra l'incertezza dei mercati e la fiducia pubblica, le stime della volatilità vengono usate dai responsabili di 
politica economica come «barometro» per misurare la vulnerabilità dei mercati e le banche centrali la monitorano per prendere decisioni 
di politica monetaria. Per questo motivo la volatilità dei mercati finanziari può avere ampie ripercussioni sull'intera economia.

Una caratteristica problematica della volatilità azionaria è che essa, a differenza dei prezzi o dei rendimenti giornalieri, non è direttamente osservabile e rimane una \gls{variabile-latente}\cite{tsay:analysis}.
Non è infatti possibile estrapolare la volatilità a partire da un singolo rendimento giornaliero.
Per ovviare a questo problema, si ricorre spesso all'uso di misure \gls{proxy}, come il valore assoluto dei rendimenti ($|r_t|$) o i rendimenti al quadrato ($r_t^2$) per ottenere una misura della varianza.
Tuttavia, la necessità di ricorrere a tali approssimazioni, data l'impossibilità di osservare il dato reale, rende complessa la valutazione delle performance dei modelli previsionali.

Sebbene non sia direttamente osservabile, essa presenta alcune caratteristiche riscontrabili nei rendimenti degli asset, come la presenza di \emph{volatility clusters} (fasi di alta volatilità alternate a periodi di calma) e una reazione asimmetrica agli shock, nota come effetto leva.

Queste caratteristiche, definibili come «fatti stilizzati», evidenziano le debolezze dei modelli tradizionali e rendono necessaria l'adozione di modelli avanzati per catturare la volatilità condizionata.


\section[Fatti stilizzati]{Fatti stilizzati delle serie storiche finanziarie}
Le proprietà statistiche che emergono regolarmente nello studio delle serie storiche, comuni a diversi strumenti, mercati e orizzonti temporali 
prendono il nome di fatti stilizzati\cite{cont:empirical}.

Conoscere tali caratteristiche è necessario per costruire modelli econometrici adeguati. I modelli finanziari tradizionali, infatti, 
si basano sull'assunzione di normalità dei rendimenti e varianza costante nel tempo (\gls{omoschedasticita}). Tuttavia, l'osservazione dei fatti 
stilizzati invalida queste ipotesi di base, rendendo i modelli lineari classici inadeguati a misurare correttamente il rischio.

Per giustificare l'adozione di modelli avanzati in grado di catturare l'evoluzione della volatilità condizionata, nelle sezioni seguenti 
vengono illustrate le principali proprietà statistiche osservate nelle serie finanziarie, con particolare attenzione alla stazionarietà, 
alla leptocurtosi, al \emph{volatility clustering} e all'effetto leva.


\subsection{Stazionarietà}
La stazionarietà è un requisito di base per l'analisi delle serie storiche: un processo è stazionario quando media, varianza e struttura di 
autocorrelazione restano costanti nel tempo. Quando queste proprietà cambiano nel tempo, la serie è non stazionaria, spesso a causa della 
presenza di \emph{trend} o radici unitarie. Tale condizione si verifica empiricamente tramite appositi test, come quello di Dickey e 
Fuller\cite{dickey:estimators} e la sua estensione, l'\gls{adf}.

Per i rendimenti finanziari, la stazionarietà garantisce l'esistenza di una varianza incondizionata finita, attorno alla quale si muove la 
volatilità condizionata. A seguito di uno shock sui prezzi, l'incertezza non aumenta senza limite, ma tende a tornare progressivamente verso 
il suo livello medio di lungo periodo.

Nei modelli GARCH che verranno trattati, Bollerslev\cite{bollerslev:garch} mostra che, con coefficienti non negativi, il processo è stazionario 
in senso lato quando la somma dei coefficienti ARCH e GARCH è minore di 1. In questo caso la varianza incondizionata è costante e finita 
e gli shock di mercato hanno un effetto solo temporaneo sulla volatilità.


\subsection{Leptocurtosi}
La curtosi è un indicatore utile per descrivere la forma di una distribuzione di rendimenti: misura il peso delle code rispetto alla normale, 
cioè la frequenza con cui si manifestano valori estremi\cite{tsay:analysis}.

Mentre una distribuzione normale standard presenta una curtosi teorica pari a 3, le serie storiche finanziarie mostrano tipicamente una 
struttura leptocurtica. Ciò si traduce in un eccesso di curtosi positivo ($K - 3 > 0$): le code risultano più pesanti e i valori estremi, 
sia positivi che negativi, occorrono con una probabilità superiore a quella prevista dal modello gaussiano\cite{cont:empirical}.

In passato si tendeva a considerare le grandi variazioni come \emph{outlier} generati da meccanismi speciali, cercando modi per escluderle 
dal campione\cite{fama:behavior}.
Questi movimenti, però, non possono essere scartati, perché possono pesare molto sul rendimento complessivo di un periodo e vanno 
modellati con attenzione, ad esempio per calcolare il \gls{var}\cite{cont:empirical}.

La natura non gaussiana dei dati rende insufficiente l'uso della deviazione standard per misurare la dispersione\cite{cont:empirical}. 
Inoltre, assumere la normalità dei rendimenti porta a sottostimare la probabilità degli eventi estremi; per questo, prima di applicare 
modelli per la volatilità come gli ARCH\cite{engle:arch}, conviene verificare la non normalità dei rendimenti.

La presenza di code pesanti, suggerita da una curtosi maggiore di 3, viene poi accertata tramite test statistici, come il test di Jarque-Bera, 
che valuta congiuntamente asimmetria e curtosi\cite{jarque:bera}, e con strumenti grafici come il QQ-plot.

I modelli GARCH risultano adatti all'analisi finanziaria perché, come i modelli ARCH, generano una distribuzione incondizionata 
con code più pesanti di quelle della normale, anche quando gli errori sono gaussiani\cite{tsay:analysis}.

% immagine
\begin{figure}[H]
    \centering
    \includegraphics[width=0.8\textwidth]{QQ_Plot_S_P_500_SPY_.png}
    \caption[QQ Plot dei rendimenti dell'indice S\&P 500]{QQ Plot dei rendimenti dell'indice S\&P 500 (rispetto alla distribuzione normale).}
    \label{fig:qqplot-sp500}
\end{figure}


\subsection{Volatility clustering}
Il \emph{volatility clustering} (volatilità a grappoli) è un fenomeno che caratterizza le serie storiche finanziarie. 
Come osservato da Mandelbrot\cite{mandelbrot:variation}, le grandi variazioni di prezzo tendono ad essere seguite da altre grandi variazioni 
mentre piccole fluttuazioni da variazioni altrettanto contenute, indipendentemente dal segno, mostrando che eventi ad alta volatilità tendono 
a raggrupparsi nel tempo.

L'assenza di autocorrelazione nei rendimenti lineari non implica assoluta indipendenza, al contrario funzioni non lineari semplici, 
come rendimenti al quadrato o in valore assoluto, mostrano un'autocorrelazione positiva. 
Questa autocorrelazione positiva prolungata è coerente con la presenza di \emph{volatility clustering}\cite{cont:empirical}.

Come già osservato nel Paragrafo \ref{sec:volatilita} (\nameref{sec:volatilita}), la volatilità giornaliera non è direttamente osservabile dai soli rendimenti di chiusura, 
il che rende complesso lo studio di questa dinamica. 
I modelli tradizionali risultano inadatti a questo scopo poiché assumono varianza condizionata costante nel tempo.

Per superare questo limite Engle\cite{engle:arch} introdusse la classe dei processi ARCH, passando al concetto di varianza condizionata 
($\sigma_t^2$), ovvero una varianza che varia nel tempo sulla base delle informazioni passate. 
Come descritto in Tsay\cite{tsay:analysis}, l'estensione ai modelli GARCH permette di catturare questo fenomeno, dato che un ampio valore 
dei rendimenti passati al quadrato genera grande varianza nel periodo successivo, replicando così la natura a grappoli della volatilità 
osservata nei mercati.


\subsection{Asimmetria ed effetto leva}
Un'altra caratteristica ricorrente delle serie storiche finanziarie è l'asimmetria della volatilità in risposta agli shock. 
Come evidenziato da Cont\cite{cont:empirical}, i rendimenti sono correlati negativamente con la variazione della volatilità futura: 
questo fenomeno è noto come effetto leva (\emph{leverage effect}) ed è stato inizialmente documentato da Black\cite{black:studies}.

Come riassunto da Nelson\cite{nelson:conditional}, la volatilità tende ad aumentare dopo cattive notizie e a diminuire dopo buone notizie. 
Nonostante l'efficacia nel catturare il fenomeno di \emph{volatility clustering}, i modelli GARCH tradizionali risultano inadatti in quanto 
assumono che variazioni positive e negative della stessa entità producano lo stesso impatto sulla varianza futura\cite{tsay:analysis}.

Per superare questo problema di simmetria e catturare l'effetto leva sono state sviluppate estensioni asimmetriche, come il GJR-GARCH e 
l'EGARCH, che vengono trattate nel capitolo successivo.